\section{Folded Gaussian blocks}
\label{sec:blocks}

We next turn the uniform Boolean sampler into a distribution adapted to
Gaussian weights. The folded Gaussian and biased-bit description follow
the setup of Li and Liu~\cite[Section~3]{LL} and
\cite[Section~2]{LLC}; all estimates required for the hierarchy are
proved here.  Throughout this section all constants, including the
lower threshold on the integer block size $b$, are absolute.  For a finite
set or group, an expectation with an unqualified subscript denotes the
uniform average.

Each point of a unit cube has two adjacent integer corrections in each
coordinate, zero and one. We attach a Gaussian weight to every resulting
bit vector. Summing these weights folds the Gaussian onto the cube; their
normalized values form a product law on the bits. We retain cube locations
where this sum has its typical size and no single bit vector carries too
much weight. After discretizing the cube, a line subgroup lets us choose
at most one usable location in each coset while retaining almost all
of the mean weight.
The preimage of that subgroup is one lattice. The next section will use
these block lattices together.

\subsection{The folded distribution and its grid}

Set
\begin{equation}
 h_b=\sqrt{\frac94\log b},\qquad
 f_b(t)=h_b\gamma_1(h_bt)+h_b\gamma_1(h_b(t-1)),\qquad
 F_b(t)=\prod_{j=1}^b f_b(t_j),
 \label{eq:folded-density}
\end{equation}
where $0\le t\le1$ in the definition of $f_b$ and
$t\in[0,1]^b$ in that of $F_b$.  At $t\in[0,1]^b$, let $P_{b,t}$ be the
product probability distribution on $\{0,1\}^b$ whose $j$th coordinate has
probabilities
\[
 P_{b,t}(\eps_j=e)
   =\frac{h_b\gamma_1(h_b(t_j-e))}{f_b(t_j)},\qquad e\in\{0,1\}.
\]
Thus the following identity holds for every bit vector:
\begin{equation}
 h_b^b\gamma_b\bigl(h_b(t-\eps)\bigr)
       =F_b(t)P_{b,t}(\eps).
 \label{eq:folded-identity}
\end{equation}
Write
\[
 a_b=\int_0^1 f_b(t)\,dt,\qquad
 \mu_b=\frac1{a_b}\int_0^1f_b(t)\log f_b(t)\,dt.
\]

\begin{lemma}[Folded Gaussian grid]
\label{lem:folded-grid}
For every sufficiently large $b$, there is a prime $p_b$ satisfying
\begin{equation}
 b\mu_b-4b^{0.56}\le\log p_b\le b\mu_b-3b^{0.56}.
 \label{eq:block-prime}
\end{equation}
Let $G_b=\F_{p_b}^b$, and associate with $z\in G_b$ the anchor
$t_z=z/p_b\in[0,1)^b$, using the representatives
$0,\ldots,p_b-1$ in every coordinate.  Define the eligible set
$E_b\subset G_b$ by the two conditions
\begin{equation}
 \left|\log F_b(t_z)-b\mu_b\right|\le b^{0.56},\qquad
 \#\left\{j:\left|t_{z,j}-\frac12\right|\le h_b^{-2}\right\}
       \ge b^{0.70}.
 \label{eq:block-eligible}
\end{equation}
Then
\begin{equation}
 \E_{z\in G_b}F_b(t_z)=1+O(b^{-1/20}),\qquad
 \E_{z\in G_b}F_b(t_z)\ind_{E_b}(z)=1+O(b^{-1/20}).
 \label{eq:block-mass}
\end{equation}
Moreover, $\mu_b=\log h_b+O(1)$.  Put
$\eta_b=bh_b^2/p_b$.  If $t$ belongs to the grid cell
$Q_z=t_z+[0,p_b^{-1})^b$, then, for every $\eps\in\{0,1\}^b$,
\begin{equation}
 \begin{split}
 e^{-\eta_b}F_b(t_z)&\le F_b(t)\le e^{\eta_b}F_b(t_z),\\
 e^{-\eta_b}\frac{F_b(t_z)}{p_b}P_{b,t_z}(\eps)
 &\le\frac{h_b^b}{p_b}\gamma_b\bigl(h_b(t-\eps)\bigr)
 \le e^{\eta_b}\frac{F_b(t_z)}{p_b}P_{b,t_z}(\eps).
 \end{split}
 \label{eq:block-rounding}
\end{equation}
There is an absolute $c>0$ such that, for every eligible $z$,
every $t\in Q_z$, and every $\eps\in\{0,1\}^b$,
\begin{equation}
 \frac{F_b(t_z)}{p_b}P_{b,t_z}(\eps)\le e^{-c b^{0.70}},
 \qquad
 \frac{h_b^b}{p_b}\gamma_b\bigl(h_b(t-\eps)\bigr)
       \le e^{-c b^{0.70}}.
 \label{eq:block-atom}
\end{equation}
\end{lemma}

\begin{proof}
First, if $X$ is a standard normal random variable, then
$a_b=\Prob(|X|\le h_b)$.  The map
\[
 X\longmapsto
 \begin{cases}
 X/h_b,&X\ge0,\\
 1+X/h_b,&X<0
 \end{cases}
\]
pushes its conditional law on $[-h_b,h_b]$ to the probability density
$f_b/a_b$ on $[0,1]$.  At the resulting point $t$, one summand in
$f_b(t)$ equals $h_b\gamma_1(X)$, whereas
$f_b(t)\le2h_b/\sqrt{2\pi}$ everywhere.  Consequently
\[
 \log h_b-\frac12\log(2\pi)
       -\frac12\E[X^2\mid |X|\le h_b]
 \le\mu_b\le\log h_b+\log\frac2{\sqrt{2\pi}}.
\]
The conditional second moment is bounded by an absolute constant, proving
$\mu_b=\log h_b+O(1)$.  The elementary Gaussian tail estimate also gives
\begin{equation}
 1-a_b\le C h_b^{-1}e^{-h_b^2/2}
       =C h_b^{-1}b^{-9/8},\qquad
 \int_{[0,1]^b}F_b(t)\,dt=a_b^b=1-O(b^{-1/8}).
 \label{eq:folded-total-mass}
\end{equation}

Here is an elementary way to choose the prime, avoiding any quantitative
prime-gap input.  Since $\mu_b\ge\log h_b-C\to\infty$,
$T=b\mu_b-3b^{0.56}$ is positive for large $b$.  Set
$l=\lfloor e^T/2\rfloor$.  For a prime $q$, the factorial valuation formula
gives
\[
 v_q\binom{2l}{l}
  =\sum_{k\ge1}\left(\left\lfloor\frac{2l}{q^k}\right\rfloor
                -2\left\lfloor\frac l{q^k}\right\rfloor\right)
  \le\lfloor\log_q(2l)\rfloor.
\]
Thus the power of each prime in the binomial coefficient is at most $2l$.
If all its prime divisors were at most $l/\log(2l)$, their number would be
at most $l/\log(2l)$, and hence
$\log\binom{2l}{l}\le l$.  This contradicts
$\binom{2l}{l}\ge4^l/(2l+1)>e^l$ for large $l$.
There is therefore a prime divisor $p_b$ with
\[
 \frac l{\log(2l)}<p_b\le2l.
\]
It follows that $T-O(\log T)\le\log p_b\le T$.  Since
$\mu_b=O(\log\log b)$, we have $\log T=O(\log b)=o(b^{0.56})$;
this proves Equation~\eqref{eq:block-prime}.  In particular $p_b$ grows
faster than every fixed power of $b$, and so does $1/\eta_b$.

To prove the mass assertions, first work in the continuous cube with
probability density $F_b/a_b^b$.  Its coordinates are independent, and
$\log F_b$ has mean $b\mu_b$.  The function $\log f_b$ takes values in an
interval of length $O(1+h_b^2)$: the upper bound was given above, and a
nearest-endpoint summand gives
$f_b(t)\ge h_b(2\pi)^{-1/2}e^{-h_b^2/8}$.
Consequently
\[
 \operatorname{Var}(\log F_b)\le Cb(\log b)^2,
 \qquad
 \Prob\left(\left|\log F_b-b\mu_b\right|>
                    \frac12b^{0.56}\right)
       \le C b^{-0.12}(\log b)^2=O(b^{-1/20}).
\]
For one coordinate, on the interval
$|t-1/2|\le(2h_b^2)^{-1}$ we have
$f_b(t)\ge c h_b e^{-h_b^2/8}$, with an absolute $c>0$.
Its probability under $f_b/a_b$ is therefore at least
$c h_b^{-1}b^{-9/32}$.  The number $N$ of coordinates in this interval is
binomial with mean
\[
 \E N\ge c h_b^{-1}b^{23/32}\ge2b^{0.70}
\]
for sufficiently large $b$, since $23/32>0.70$.  Its variance is at most
its mean, so
\[
 \Prob(N<b^{0.70})
       \le\frac4{\E N}
       \le C h_b b^{-23/32}=O(b^{-1/20}).
\]
Thus a set of continuous $F_b$-mass $1-O(b^{-1/20})$ satisfies the
logarithmic eligibility condition with half its tolerance and the
balanced-coordinate condition with half its interval width.

For each $e\in\{0,1\}$ and $u\in[0,1]$,
\[
 \left|\frac{d}{du}\log\bigl(h_b\gamma_1(h_b(u-e))\bigr)\right|
        =h_b^2|u-e|\le h_b^2.
\]
The same bound holds for $|(\log f_b)'|$, since its derivative is a
weighted average of these two logarithmic derivatives.  Integrating
along the coordinates of a grid cell proves
Equation~\eqref{eq:block-rounding}.  Rounding a continuous point down to
its anchor changes $\log F_b$ by at most $\eta_b$ and every coordinate by
less than $p_b^{-1}$.  Since
$\eta_b<\frac12b^{0.56}$ and $p_b^{-1}<(2h_b^2)^{-1}$ for large $b$,
the stricter continuous event just considered rounds into $E_b$.
In addition, summing the first bound in
Equation~\eqref{eq:block-rounding} over all cells gives
\[
 e^{-\eta_b}\E_zF_b(t_z)
 \le\int_{[0,1]^b}F_b(t)\,dt
 \le e^{\eta_b}\E_zF_b(t_z).
\]
Applying the corresponding comparison on cells whose anchors lie in
$E_b$ gives the lower bound on their weighted grid mass.  Its upper bound
is the total grid mass.  Equation~\eqref{eq:folded-total-mass} now proves
both assertions in Equation~\eqref{eq:block-mass}.

Finally, the odds of a coordinate bit at $t$ are
\begin{equation}
 \frac{P_{b,t}(\eps_j=0)}{P_{b,t}(\eps_j=1)}
       =\exp\bigl(h_b^2(1/2-t_j)\bigr).
 \label{eq:block-bit-odds}
\end{equation}
On a balanced coordinate at an eligible anchor, each bit value has
probability at least $\alpha=(1+e)^{-1}$.  Therefore
$P_{b,t_z}(\eps)\le(1-\alpha)^{b^{0.70}}$ for every $\eps$.
Equations~\eqref{eq:block-prime} and~\eqref{eq:block-eligible} give
$F_b(t_z)/p_b\le e^{5b^{0.56}}$.  Since $0.56<0.70$, their product is
at most $e^{-c b^{0.70}}$ after decreasing an absolute $c>0$ and increasing
the absolute threshold.  Equation~\eqref{eq:block-rounding} proves the
same assertion at all points of the cell, again decreasing $c$ if needed.
\end{proof}

The exponents are chosen to leave room between four different estimates:
\[
 \frac12<0.56,\qquad
 \frac{0.56}{0.9}<0.65<0.69<0.70<\frac{23}{32}.
\]
The information window exceeds the square-root fluctuation scale; the
previous block's weight cap fits the sampler; and the number of balanced
coordinates exceeds the number of fair bits needed.  The first block will
have size about $(\log\log n)^2$, so $2(0.70)>1$ also makes its atom
bound stronger than the final logarithmic loss.  None of these values
is intended to be optimal.

\subsection{Usable anchors and sparse line maxima}

The next lemma has two roles.  Most of the folded Gaussian mass is kept on
anchors whose biased bits sample a prescribed preceding block uniformly
in its base point.  A suitable line in the new grid then permits this mass
to be represented by one selected anchor in each coset.

\begin{lemma}[Usable anchors]
\label{lem:usable}
There are absolute constants $A,B>0$ with the following property.
Let $b\ge B$ be an integer, and use the objects in
Lemma~\ref{lem:folded-grid}.  Suppose that $H$ is a finite abelian group,
$\log|H|\le b^4$, and $g_0:H\to[0,\infty)$ satisfies
\[
 \frac12\le\bar g_0:=\E_{x\in H}g_0(x)\le2,\qquad
 g_0(x)\le\exp(5a^{0.56})\quad(x\in H)
\]
for some $1\le a\le b^{1/0.9}$.  Then one can choose shifts
$w_1,\ldots,w_b\in H$, a usable set $U_b\subset E_b$, and a
one-dimensional linear subspace $C_b\subset G_b$ such that, writing
$W\eps=\sum_{j=1}^b\eps_jw_j$,
\begin{equation}
 \sum_{\eps\in\{0,1\}^b}P_{b,t_z}(\eps)g_0(x-W\eps)
       \ge(1-A b^{-1/100})\bar g_0
       \qquad(z\in U_b,\ x\in H).
 \label{eq:block-biased-sampler}
\end{equation}
Define
\begin{equation}
 v_b(z)=\frac{F_b(t_z)}{p_b}\ind_{U_b}(z),\qquad
 g_b(x)=\max_{c\in C_b}v_b(x-c),\qquad
 \bar g_b=\E_{x\in G_b}g_b(x).
 \label{eq:block-line-weight}
\end{equation}
These choices satisfy
\begin{equation}
 |\bar g_b-1|\le A b^{-1/100},\qquad
 0\le g_b(x)\le\exp(5b^{0.56})\quad(x\in G_b),\qquad
 \left|\E_z F_b(t_z)\ind_{U_b}(z)-1\right|
       \le A b^{-1/20}.
 \label{eq:block-usable-bounds}
\end{equation}
There is also a base version with no $H$, $g_0$, or $W$: take $U_b=E_b$
and choose $C_b$ so that the conclusions concerning $g_b$ in
Equation~\eqref{eq:block-usable-bounds} hold.
\end{lemma}

\begin{proof}
We first select $W$ and $U_b$ when preceding data are given.
For large $b$, the bound on $a$ implies
\[
 5a^{0.56}\le5b^{0.56/0.9}\le b^{0.65},
\]
because $0.56/0.9<0.65$.  Thus $g_0$ satisfies the cap in
Lemma~\ref{lem:sampler}.  Choose the columns $w_j$ independently and
uniformly from $H$.

Fix an eligible anchor $z$.  A bit whose probability of being $1$ is $q$
has the following exact mixture representation: with probability
$2\min(q,1-q)$ use a fair bit, and otherwise use its more likely value.
Make the mixture choices independently at the $b$ coordinates.  Write
$S\subset\{1,\ldots,b\}$ for the coordinates chosen to be fair, and
$d_j$ for the deterministic value at $j\notin S$.
Conditional on these choices, the resulting vector has fair independent
bits on $S$ and the values $d_j$ elsewhere; averaging these conditional
laws recovers $P_{b,t_z}$ exactly.

By Equation~\eqref{eq:block-bit-odds}, at least $b^{0.70}$ coordinates
are selected into $S$ with probability at least $2\alpha$, where
$\alpha=(1+e)^{-1}$.  Apply Chebyshev's inequality to the number selected
among these coordinates.  Its mean is at least $2\alpha b^{0.70}$,
its variance is at most its mean, and $b^{0.69}$ is at most half its mean
for sufficiently large $b$.  Hence
\begin{equation}
 \Prob\bigl(|S|<b^{0.69}\bigr)=O(b^{-0.70})=O(b^{-1/2}).
 \label{eq:block-fair-count}
\end{equation}
For any fixed $S$ with $b^{0.69}\le|S|\le b$, its columns remain
independent uniform elements of $H$.  Lemma~\ref{lem:sampler} therefore
shows, with failure probability $O(b^{-1/2})$ over $W$, that
\[
 2^{-|S|}\sum_{\eps\in\{0,1\}^S}
      g_0\left(y-\sum_{j\in S}\eps_jw_j\right)
       \ge(1-Cb^{-1/10})\bar g_0\qquad(y\in H).
\]
The conclusion is simultaneous for every $y$.  In particular it allows
$y=x-\sum_{j\notin S}d_jw_j$, even though that offset involves the other
columns of $W$.  There is no conditioning on those columns in applying
the sampler to the selected ones.

Let $q_W(z)$ be the mixture probability that either $|S|<b^{0.69}$ or
the displayed uniform sampler conclusion fails for its selected columns.
Equation~\eqref{eq:block-fair-count} and Lemma~\ref{lem:sampler} imply
$\E_W q_W(z)\le Cb^{-1/2}$ for each eligible $z$.  The weighted average
is therefore bounded by
\[
 \E_W\E_{z\in G_b}
       F_b(t_z)\ind_{E_b}(z)q_W(z)
       \le Cb^{-1/2}\E_zF_b(t_z)\ind_{E_b}(z)
       \le C' b^{-1/2}.
\]
Fix one $W$ for which the inner expectation is at most this last bound,
and set
\[
 U_b=\{z\in E_b:q_W(z)\le b^{-1/20}\}.
\]
Weighted Markov's inequality gives the explicit discarded mass
\begin{equation}
 \E_zF_b(t_z)\ind_{E_b\setminus U_b}(z)
       \le C' b^{-1/2+1/20}=C'b^{-9/20}.
 \label{eq:block-discarded-mass}
\end{equation}
Together with Equation~\eqref{eq:block-mass}, this proves the last
assertion of Equation~\eqref{eq:block-usable-bounds}.  At every remaining
anchor, average the uniform sampler conclusion over successful mixture
components and use nonnegativity on the others.  For all $x\in H$ the
result is at least
\[
 (1-b^{-1/20})(1-Cb^{-1/10})\bar g_0
       \ge(1-A b^{-1/100})\bar g_0,
\]
proving Equation~\eqref{eq:block-biased-sampler}.

It remains to choose $C_b$.  The following argument applies both to this
$U_b$ and, in the base version, to $U_b=E_b$.  Put
\[
 M=\E_zF_b(t_z)\ind_{U_b}(z)=1+O(b^{-1/20}),\qquad
 d=\frac{|U_b|}{p_b^b}.
\]
The lower eligible weight gives
\begin{equation}
 d\le M e^{-b\mu_b+b^{0.56}}
       \le2e^{-b\mu_b+b^{0.56}}.
 \label{eq:block-sparse-density}
\end{equation}
Choose a line $C\subset G_b$ uniformly among one-dimensional linear
subspaces.  For any fixed nonzero $u\in G_b$,
\[
 \Prob_C(u\in C)=\frac{p_b-1}{p_b^b-1}.
\]
Consequently, for every $z\in U_b$, the event $\mathcal C_C(z)$ that
$z+C$ contains another point of $U_b$ has probability at most
\begin{equation}
 \Prob_C(\mathcal C_C(z))
       \le(|U_b|-1)\frac{p_b-1}{p_b^b-1}
       \le2p_b d\le4e^{-2b^{0.56}}.
 \label{eq:block-collision-probability}
\end{equation}
The last inequality uses the upper bound in
Equation~\eqref{eq:block-prime}.

For a fixed line, write
$S_C(x)=\sum_{c\in C}v_b(x-c)$ and
$g_C(x)=\max_{c\in C}v_b(x-c)$.
The normalization in Equation~\eqref{eq:block-line-weight} ensures
\[
 \E_x S_C(x)=p_b\E_z v_b(z)=M.
\]
If a coset contains at most one usable anchor, its sum and maximum agree.
Otherwise, charge their difference to all usable weights in that coset.
This gives the pointwise bound
\[
 0\le S_C(x)-g_C(x)
       \le\sum_{c\in C}v_b(x-c)\ind_{\mathcal C_C(x-c)}.
\]
Averaging first over $x$ and then over $C$ yields
\begin{align}
 \E_C\E_x(S_C(x)-g_C(x))
 &\le\E_C\left[p_b\E_z v_b(z)\ind_{\mathcal C_C(z)}\right]\notag\\
 &=\E_zF_b(t_z)\ind_{U_b}(z)\Prob_C(\mathcal C_C(z))\notag\\
 &\le4M e^{-2b^{0.56}}.
 \label{eq:block-weighted-collision}
\end{align}
In particular the $p_b^{-1}$ in $v_b$ cancels the $p_b$ translates in the
coset sum; the estimate controls lost weight, not just the number of
colliding cosets.  Some line $C_b$ thus satisfies
\[
 M-4M e^{-2b^{0.56}}\le\bar g_b\le M.
\]
This proves the mean assertion, in fact with the stronger error
$O(b^{-1/20})$.  Finally every usable anchor is eligible, so
$v_b(z)\le e^{5b^{0.56}}$ by
Equations~\eqref{eq:block-prime} and~\eqref{eq:block-eligible}; the same
bound holds for its coset maximum.  Enlarging the absolute constants
$A,B$ completes both versions of the lemma.
\end{proof}

\begin{remark}
The function $g_b$ need not have a positive lower bound at every point:
many cosets can contain no usable anchor.  Its mean is close to one, and
Equation~\eqref{eq:block-biased-sampler} is what recovers that mean
uniformly at the next stage.  This distinction is essential when the
blocks are assembled.
\end{remark}

For use in that assembly, the grid construction has the following exact
interpretation in a lattice.  We identify elements of $C_b$ with their
integer representatives when writing $C_b/p_b$.

\begin{lemma}[From a coset to a Gaussian cell]
\label{lem:block-cell}
Let $L_b=\Z^b+C_b/p_b$.  For $t\in\R^b$ put
$k=\lfloor p_b t\rfloor\in\Z^b$ and $x=k\pmod{p_b}\in G_b$.
Floors and fractional parts are taken coordinatewise.
If $g_b(x)>0$, choose $c\in C_b$ attaining its maximum and let
$z=x-c\in U_b$, again using standard representatives.  Then
\begin{equation}
 \lambda=\frac{k-z}{p_b}\in L_b,\qquad
 \tau=t-\lambda=\frac{z+\{p_b t\}}{p_b}\in Q_z,\qquad
 g_b(x)=\frac{F_b(t_z)}{p_b}.
 \label{eq:block-cell-identity}
\end{equation}
For all $\eps\in\{0,1\}^b$, the point $\lambda+\eps$ lies in $L_b$,
its residual $t-\lambda-\eps=\tau-\eps$ has coordinates in $[-1,1]$,
and
\begin{equation}
 \frac{h_b^b}{p_b}\gamma_b\bigl(h_b(t-\lambda-\eps)\bigr)
       \ge e^{-\eta_b}g_b(x)P_{b,t_z}(\eps).
 \label{eq:block-cell-lower}
\end{equation}
The corresponding individual weight is at most $e^{-c b^{0.70}}$.
\end{lemma}

\begin{proof}
The integer vector $k-z$ reduces to $c$ modulo $p_b$, proving
$\lambda\in L_b$.  The remaining identities in
Equation~\eqref{eq:block-cell-identity} follow from
$p_b t=k+\{p_b t\}$ and the chosen maximizing anchor.  Adding an integer
bit vector preserves $L_b$.  Since $\tau\in[0,1)^b$, the residual
coordinate bound follows.  Apply
Equations~\eqref{eq:block-rounding} and~\eqref{eq:block-atom} at
$\tau\in Q_z$ to obtain the stated lower and upper estimates.
\end{proof}

In the hierarchy below, the preceding group is $H=G_a$ with
$a\le b^{1/0.9}$.  Its size satisfies the remaining hypothesis of
Lemma~\ref{lem:usable}, because
\[
 \log|G_a|=a\log p_a=O(a^2\log\log a)
       \le O(b^{2/0.9}\log\log b)\le b^4
\]
for sufficiently large $b$.  Both this bound and the cap comparison in
the proof have strict exponent margins.  Thus the same absolute cutoff
works at every level, independently of the number of blocks.
